\section{Cyclotomic cosets modulo powers of ten}\label{sec:cosets-ten}

For $q$ coprime to $10$, let $C_q(n)$ be the number of distinct
orbits of the permutation $x\mapsto qx$ on $\mathbb Z/10^n\mathbb Z$.
Thus the orbit containing $x$ is the finite set of all $q^kx$ with
$k\ge0$. The zero orbit is included. V.~Raman's sequences
A220018, A220019, A220020, and A221855 correspond respectively to
$q=3,7,9,13$. Colin Barker recorded the following conjectures on
13--14 April 2013 \cite{oeis220018,oeis220019,oeis220020,oeis221855}.

\begin{theorem}\label{thm:cosets-ten}
For every $n\ge3$,
\[
 C_3(n)=8n^2-14n-1,\qquad
 C_9(n)=4(4n^2-7n-1),\qquad
 C_{13}(n)=2n(4n-7).
\]
For every $n\ge4$,
\[
 C_7(n)=5(16n^2-60n+37).
\]
\end{theorem}

The ingredients are standard orbit counting, the Chinese remainder
theorem, and lifting the exponent. A much more general classification
of cyclotomic cosets and their sizes appears in Zhu--Liu--Wu
\cite{zhu-liu-wu-cosets}. The present argument gives the explicit
specialization required by these four recorded formulas, with a
complete Lean proof from the literal orbit definition.

\begin{lemma}[Divisor formula]\label{lem:cosets-divisors}
For positive $m$ coprime to $q$, the number of orbits of multiplication
by $q$ on $\mathbb Z/m\mathbb Z$ is
\[
 \sum_{d\mid m}\frac{\varphi(d)}{\operatorname{ord}_d(q)}.
\]
Here $\varphi(1)=\operatorname{ord}_1(q)=1$.
\end{lemma}
\begin{proof}
The additive cyclic group $\mathbb Z/m\mathbb Z$ has exactly
$\varphi(d)$ elements of additive order $d$ for each divisor $d$ of
$m$. Multiplication by a unit preserves additive order. If $x$ has
additive order $d$, then
\[
 q^kx=x\quad\Longleftrightarrow\quad(q^k-1)x=0
 \quad\Longleftrightarrow\quad d\mid q^k-1.
\]
The least positive such $k$ is $\operatorname{ord}_d(q)$.
The $\varphi(d)$ elements in the stratum therefore partition into
orbits of that common size, giving the displayed quotient.
Summing over all divisors counts every orbit exactly once.
The stratum $d=1$ consists only of zero.
\end{proof}

Write $t_+=\max(t,0)$. The exact orders needed in the proof are
\begin{equation}\label{eq:cosets-orders-two}
\begin{array}{c|c|c}
q& i=0,1&i\ge2\\\hline
3&1&2\cdot2^{(i-3)_+}\\
7&1&2\cdot2^{(i-4)_+}\\
9&1&2^{(i-3)_+}\\
13&1&2^{(i-2)_+}
\end{array}
\end{equation}
and, for $j\ge1$,
\begin{equation}\label{eq:cosets-orders-five}
\begin{array}{c|c}
q&\operatorname{ord}_{5^j}(q)\\\hline
3&4\cdot5^{j-1}\\
7&4\cdot5^{(j-2)_+}\\
9&2\cdot5^{j-1}\\
13&4\cdot5^{j-1}.
\end{array}
\end{equation}
The orders at $j=0$ are $1$.

\begin{proof}[Proof of the order formulas]
The orders of $3,7$ modulo $4$ are $2$, and their exact valuations are
$v_2(3^2-1)=3$ and $v_2(7^2-1)=4$. Lifting the exponent gives
\[
 v_2(3^{2k}-1)=3+v_2(k),\qquad
 v_2(7^{2k}-1)=4+v_2(k).
\]
For $i\ge2$, an exponent producing $1$ modulo $2^i$ must be even.
The displayed laws therefore show that the least possible exponent is
$2\cdot2^{(i-3)_+}$ or $2\cdot2^{(i-4)_+}$, respectively.
For $9$ and $13$, which are $1$ modulo $4$, the analogous laws are
\[
 v_2(9^k-1)=3+v_2(k),\qquad
 v_2(13^k-1)=2+v_2(k).
\]
They give the other two rows, with the small moduli verified directly.

Modulo $5$, the orders of $3,7,9,13$ are $4,4,2,4$, respectively.
The exact valuations of $q^e-1$ at these exponents are $1,2,1,1$.
In particular $7^4-1=2400=5^2\cdot96$ is divisible by $25$ but
not $125$. For each of the four pairs $(q,e)$, lifting the exponent
gives
\[
 v_5(q^{ek}-1)=v_5(q^e-1)+v_5(k).
\]
Any exponent producing $1$ modulo $5^j$ is a multiple of $e$;
the smallest possible multiplier is the power of $5$ displayed in
\eqref{eq:cosets-orders-five}. These statements apply to every
positive exponent, not only a finite list of moduli.
\end{proof}

Every divisor of $10^n$ has a unique representation $2^i5^j$ with
$0\le i,j\le n$. By the Chinese remainder theorem,
\[
 \operatorname{ord}_{2^i5^j}(q)
 =\operatorname{lcm}\bigl(\operatorname{ord}_{2^i}(q),
                         \operatorname{ord}_{5^j}(q)\bigr).
\]
Put
\[
 T_q(i,j)=\frac{\varphi(2^i5^j)}{\operatorname{ord}_{2^i5^j}(q)}.
\]
Then Lemma~\ref{lem:cosets-divisors} becomes
\begin{equation}\label{eq:cosets-grid}
 C_q(n)=\sum_{i=0}^n\sum_{j=0}^n T_q(i,j).
\end{equation}
Multiplicativity of $\varphi$, together with
$\varphi(2^{i+1})=2^i$ and $\varphi(5^{j+1})=4\cdot5^j$,
gives the following complete table:
\[
\begin{array}{c|rrrrrr}
 &i=0&i=1&i=2&i=3&i=4&i\ge5\\\hline
T_3(i,0)&1&1&1&2&2&2\\
T_3(i,j),\ j\ge1&1&1&2&4&8&8\\
T_7(i,0)&1&1&1&2&4&4\\
T_7(i,1)&1&1&2&4&8&16\\
T_7(i,j),\ j\ge2&5&5&10&20&40&80\\
T_9(i,0)&1&1&2&4&4&4\\
T_9(i,j),\ j\ge1&2&2&4&8&16&16\\
T_{13}(i,0)&1&1&2&2&2&2\\
T_{13}(i,j),\ j\ge1&1&1&2&4&8&8
\end{array}
\]
For example, if $i=r+4$ and $j=s+1$, the orders for base $3$ are
$4\cdot2^r$ and $4\cdot5^s$. Since $2^r$ and $5^s$ are coprime,
their least common multiple is $4\cdot2^r5^s$. Dividing
$2^{r+3}\cdot4\cdot5^s$ by it gives $8$. The other infinite
tails follow by the same exact cancellations.

\begin{proof}[Proof of Theorem~\ref{thm:cosets-ten}]
The sum through index $n$ of $1,1,2,4,8,8,\ldots$ is
$8n-16$ for $n\ge3$. The sum through index $n$ of
$1,1,2,4,8,16,16,\ldots$ is $16n-48$ for $n\ge4$.
The four $j=0$ rows sum respectively to $2n-1$, $4n-7$,
$4n-4$, and $2n$, in the ranges required by the theorem.
Consequently \eqref{eq:cosets-grid} gives
\begin{align*}
 C_3(n)&=(2n-1)+n(8n-16)=8n^2-14n-1,\\
 C_7(n)&=(4n-7)+(1+5(n-1))(16n-48)
        =80n^2-300n+185,\\
 C_9(n)&=(4n-4)+2n(8n-16)=16n^2-28n-4,\\
 C_{13}(n)&=2n+n(8n-16)=8n^2-14n.
\end{align*}
These are exactly the four conjectured polynomials.
\end{proof}

\paragraph{Formal scope and attribution.}
The Lean definition counts the quotient of $\mathbb Z/10^n\mathbb Z$
by the action of the subgroup generated by the unit $q$.
The theorem \path{mem_coset_iff} proves that membership in these
orbits is equivalent to $y=q^kx$ for some natural $k$.
The divisor formula is then proved using orbit sizes and additive
orders. Subsequent modules prove the divisor-to-exponent bijection,
all prime power orders by exact valuations, every entry and tail of
the table, and all four sums. Thus the polynomial is the conclusion,
not the definition of any formal source sequence.

The four endpoint theorems in \path{Proof.lean} use exactly the
source ranges $n>2$ or $n>3$ and have only the standard Lean logical
axioms as dependencies. All six modules are explicit build roots.
The source entries and earlier generators already contain the
divisor algorithm or polynomial patterns; the general methods are
classical and are not counted as new results. This package contains
four related source-assertion proofs in one family. The separately
recorded bounded prior-work search found no earlier explicit
settlement of these four formulas and does not establish worldwide
priority.
